\labsection{4}{Mobile banking heuristic audit and funnel analysis}{sec:lab04-mobile}

\begin{labproject}{Score two banking apps against the chapter's five principles, then find where users are lost}
\textbf{Coverage.} Chapter~\ref{ch:mobile-banking}.\\
\textbf{Goal.} Convert the five usability principles into a repeatable scored instrument, and locate the single step where an onboarding journey loses most of its users.

\courselabflow{Define tasks}{Record events}{Build funnel}{Explain friction}



\end{labproject}

\textbf{Complete worked example.} This runnable fixture uses synthetic inputs. First reproduce it, then adapt the data or ratings with documented evidence. All code below belongs to one file.
\begin{lstlisting}[style=utpterminal]
python all_experiments/lab04_mobile_banking_funnel.py
\end{lstlisting}

\noindent\textbf{Setup and output folder.}
\lstinputlisting[style=utppython,firstline=1,lastline=10]{all_experiments/lab04_mobile_banking_funnel.py}

\noindent\textbf{Generate one ordered cohort.}
\lstinputlisting[style=utppython,firstline=11,lastline=25]{all_experiments/lab04_mobile_banking_funnel.py}

\noindent\textbf{Validate sequence before counting conversion.}
\courselabtrace{Count unique users at every stage}{Entry}{Use a fixed cohort and one observation window}{Progress}{Count each user once per completed stage}{Drop-off}{Previous-stage count minus current-stage count}



\textbf{Use the chapter while coding.} Revisit \hyperref[ch:mobile-banking]{the earlier chapter explanation}; match its operation to the variables and state changes below.

\textbf{Follow the supplied implementation.} The sequence validation rejects missing earlier steps before rates are calculated. The first step has no previous denominator, so its step conversion is missing rather than zero. Keep user counts separate from raw event-row counts.
\courselabtrace{Trace the actual cohort intersections}{1. Input / state}{cohort is the set of users at app\_open; previous initially points to that set.}{2. Operation}{For each step, reached intersects users at that step with previous; counts records its size.}{3. Check}{overall divides by the initial cohort; step\_conversion divides by the preceding count.}


\lstinputlisting[style=utppython,firstline=26,lastline=42]{all_experiments/lab04_mobile_banking_funnel.py}

\noindent\textbf{Inspect a suspected friction point using a denominator.}
\lstinputlisting[style=utppython,firstline=43,lastline=48]{all_experiments/lab04_mobile_banking_funnel.py}
\begin{labvisual}{Synthetic cohort: where users leave the funnel}
\begin{labgraphic}\begin{tikzpicture}[font=\small]
\node[anchor=east,font=\scriptsize] at (-.2,0.0) {app open};
\fill[UTPPaleBlue] (0,-0.16) rectangle (7,0.16);
\fill[UTPTeal] (0,-0.16) rectangle (7.0000,0.16);
\node[anchor=west,font=\scriptsize] at (7.15,0.0) {1000.0 users};
\node[anchor=east,font=\scriptsize] at (-.2,-0.55) {signup start};
\fill[UTPPaleBlue] (0,-0.7100000000000001) rectangle (7,-0.39);
\fill[UTPTeal] (0,-0.7100000000000001) rectangle (5.2150,-0.39);
\node[anchor=west,font=\scriptsize] at (7.15,-0.55) {745.0 users};
\node[anchor=east,font=\scriptsize] at (-.2,-1.1) {id upload};
\fill[UTPPaleBlue] (0,-1.26) rectangle (7,-0.9400000000000001);
\fill[UTPTeal] (0,-1.26) rectangle (4.4100,-0.9400000000000001);
\node[anchor=west,font=\scriptsize] at (7.15,-1.1) {630.0 users};
\node[anchor=east,font=\scriptsize] at (-.2,-1.6500000000000001) {id verified};
\fill[UTPPaleBlue] (0,-1.81) rectangle (7,-1.4900000000000002);
\fill[UTPTeal] (0,-1.81) rectangle (2.7230,-1.4900000000000002);
\node[anchor=west,font=\scriptsize] at (7.15,-1.6500000000000001) {389.0 users};
\node[anchor=east,font=\scriptsize] at (-.2,-2.2) {address entered};
\fill[UTPPaleBlue] (0,-2.3600000000000003) rectangle (7,-2.04);
\fill[UTPTeal] (0,-2.3600000000000003) rectangle (2.4570,-2.04);
\node[anchor=west,font=\scriptsize] at (7.15,-2.2) {351.0 users};
\node[anchor=east,font=\scriptsize] at (-.2,-2.75) {account funded};
\fill[UTPPaleBlue] (0,-2.91) rectangle (7,-2.59);
\fill[UTPTeal] (0,-2.91) rectangle (2.0300,-2.59);
\node[anchor=west,font=\scriptsize] at (7.15,-2.75) {290.0 users};
\node[anchor=east,font=\scriptsize] at (-.2,-3.3000000000000003) {first payment};
\fill[UTPPaleBlue] (0,-3.4600000000000004) rectangle (7,-3.14);
\fill[UTPTeal] (0,-3.4600000000000004) rectangle (1.8410,-3.14);
\node[anchor=west,font=\scriptsize] at (7.15,-3.3000000000000003) {263.0 users};
\end{tikzpicture}\end{labgraphic}\end{labvisual}

\textbf{Check the run.} The following output was captured from the supplied script; it describes the synthetic fixture only.
\lstinputlisting[style=utpterminal]{labs/outputs/lab04.txt}
% Reused from chapters/04-mobile-and-internet-banking.tex
\begin{labvisual}{Recall the chapter model}
\begin{labgraphic}
\begin{tikzpicture}[
  b/.style={draw=UTPBlue,fill=UTPPaleBlue,rounded corners=2mm,text width=2.6cm,minimum height=1.1cm,align=center,font=\scriptsize\bfseries\color{UTPNavy}},
  arr/.style={-{Stealth[length=2mm]},thick,draw=UTPTeal}
]
\node[b] (a) at (-4.5,0) {Reactive\\``show my balance''};
\node[b] (b) at (-1.5,0) {Contextual\\``show my spending''};
\node[b] (c) at (1.5,0) {Predictive\\``warn me early''};
\node[b] (d) at (4.5,0) {Assistive\\``help me act''};
\draw[arr](a)--(b); \draw[arr](b)--(c); \draw[arr](c)--(d);
\end{tikzpicture}
\end{labgraphic}\end{labvisual}


\begin{verification}
\textbf{Explain and extend your result.} Recalculate one stage from the underlying events. Identify the largest absolute drop and propose a measurable change, including the metric and denominator needed to evaluate it.

\begin{itemize}\item The script runs from a fresh process and creates the named CSV or report files.\item Every rate uses a stated population and denominator; every model assumption is identified.\item Synthetic outputs are labelled and are not presented as observations of a real institution.\item Change one input, rerun, and explain the direction of the result.\end{itemize}\end{verification}
\begin{labdeliverable}Submit the complete script, generated evidence, and a one-page interpretation linking the results to the chapter model. State the assumptions and one limitation.\end{labdeliverable}
